\documentclass[../../../../main.tex]{subfiles}
\begin{document}

\begin{flushright}
\footnotesize\textit{【自動文字起こし・要確認】}
\end{flushright}

\begin{proof}[解]
\[
\begin{array}{ccc}
k-1 & \xrightarrow{1-\frac{k-1}{n}} & k \\
r-1 & & r \\
r & \xrightarrow{\frac{r}{n}} & r
\end{array}
\]

\textbf{(1)} 上図から
\[
P(S_k=r)=\frac{n+1-r}{n}P(S_{k-1}=r-1)+\frac{r}{n}P(S_{k-1}=r) \quad(r=1,2,\cdots,k\text{の時}, P(S_0=r)=0\text{と扱う})
\]

\textbf{(2)} 確率変数$X_t$を，
\[
X_t=
\begin{cases}
0 & (t\text{のカードを}k\text{回目まで引かなかった}) \\
1 & (t\text{のカードを}k\text{回目までに引いた})
\end{cases}
\]
と定めると，$S_k=\displaystyle\sum_{t=1}^nX_t$だから，期待値の加法性から
\[
E(S_k)=E\left(\sum_{t=1}^nX_t\right)=\sum_{t=1}^nE(X_t)=nE(X_1) \quad(\because\text{対称性}) \quad\cdots\text{①}
\]
である。$k$回目までに1のカードを引く確率$\ell_k$は，排反をかんがえて，
\[
\ell_k=1-\left(\frac{n-1}{n}\right)^k
\]
で，$E(X_1)=0\cdot P(X_1=0)+1\cdot P(X_1=1)=P(X_1=1)=\ell_k$だから，①より，
\[
E_k=n\left[1-\left(\frac{n-1}{n}\right)^k\right]
\]
\end{proof}

\bigskip
\noindent\textbf{[別解]}

\begin{proof}[別解]
\textbf{(2)} (1)の両辺$r$をかけて$r$について和をとる。
\[
\sum_{r=1}^krP(S_k=r)=\sum_{r=1}^k\frac{r^2}{n}P(S_{k-1}=r)+\sum_{r=1}^k\frac{nr+r(r-1)}{n}P(S_{k-1}=r-1)
\]
$P(S_{k-1}=k)=0$だから，
\[
E_k=\sum_{r=1}^{k-1}\frac{r^2}{n}P(S_{k-1}=r)+\sum_{r=1}^k\frac{(n-1)(r-1)+n-(r-1)^2}{n}P(S_{k-1}=r-1)
\]
\[
=\sum_{r=1}^{k-1}\frac{r^2}{n}P(S_{k-1}=r)+\frac{n-1}{n}E_{k-1}+\sum_{r=1}^kP(S_{k-1}=r-1)-\sum_{r=1}^k\frac{(r-1)^2}{n}P(S_{k-1}=r-1)
\]
最初の和と最後の和は（添字を$j=r-1$とおきかえると）互いに打ち消しあうので，$\displaystyle\sum_{r=1}^kP(S_{k-1}=r-1)=1$から，
\[
E_k=\frac{n-1}{n}E_{k-1}+1
\]
だから，$E_1=1$より
\[
E_k=n\left[1-\left(\frac{n-1}{n}\right)^k\right]
\]
\end{proof}

\end{document}
